\documentclass[8pt,fleqn]{extarticle}

\usepackage[T1]{fontenc}
\usepackage[a4paper,margin=0.6cm]{geometry}
\usepackage{multicol}
\usepackage{amsmath,amssymb}
\usepackage{enumitem}
\usepackage{titlesec}
\usepackage[protrusion=true,expansion=false]{microtype}

% ---- EXTREME COMPRESSION v2 ----
\setlength{\parindent}{0pt}
\setlength{\parskip}{0pt}
\setlength{\columnsep}{1.5pt}
\setlength{\multicolsep}{0pt}
\setlength{\premulticols}{0pt}
\setlength{\postmulticols}{0pt}
\setlength{\mathindent}{0pt}

\setlength{\abovedisplayskip}{0pt}
\setlength{\belowdisplayskip}{0pt}
\setlength{\abovedisplayshortskip}{0pt}
\setlength{\belowdisplayshortskip}{0pt}

\setlength{\topsep}{0pt}
\setlength{\partopsep}{0pt}
\setlength{\itemsep}{0pt}
\setlength{\parsep}{0pt}

\renewcommand{\baselinestretch}{0.9}
\emergencystretch=0.8em

\titlespacing*{\section}{0pt}{0.2pt}{0.25em}
\titleformat{\section}[runin]{\bfseries}{}{0em}{}[.\enspace]
\raggedcolumns


\begin{document}
\tiny
\begin{multicols}{3}

% =========================
% GROUP 1: REGRESSION
% =========================
\section*{Linear Regression}
Model: $y=X\beta+\epsilon$, $\epsilon\sim\mathcal N(0,\sigma^2)$ i.i.d.
Loss: $J=\frac1{2n}\|y-X\beta\|^2$.
Closed form: $\hat\beta=(X^TX)^{-1}X^Ty$ (requires $X^TX$ invertible).
Grad: $\nabla J=-\frac1nX^T(y-X\beta)=0\Rightarrow$ normal equations $X^TX\beta=X^Ty$.
MLE $\equiv$ OLS under Gaussian noise.
Predictions: $\hat y=H y$, $H=X(X^TX)^{-1}X^T$ (hat matrix), $\text{tr}(H)=p$.

% =========================
\section*{Simple Regression Closed Form}

$\hat\beta_1=\frac{\sum(x_i-\bar x)(y_i-\bar y)}{\sum(x_i-\bar x)^2}$,\quad $\hat\beta_0=\bar y-\hat\beta_1\bar x$.

Correlation: $r=\frac{\sum(x_i-\bar x)(y_i-\bar y)}{\sqrt{\sum(x_i-\bar x)^2\sum(y_i-\bar y)^2}}$.

OLS: minimum residual sum of squares.


% =========================
\section*{Logistic Regression}
$\sigma(z)=\frac1{1+e^{-z}}$, $\sigma'=\sigma(1-\sigma)$, $z=w^Tx$.
$P(y=1|x)=\sigma(w^Tx)$; decision boundary: $w^Tx=0$ (linear).
Loss (NLL): $J=-\frac1n\sum[y\log p+(1-y)\log(1-p)]$, $p_i=\sigma(w^Tx_i)$.
Grad: $\nabla J=\frac1nX^T(p-y)$. Hessian: $X^TWX\succeq0$ $\Rightarrow$ convex.
$W=\text{diag}(p_i(1-p_i))$. No closed form; use GD/Newton's method.
Logit: $\log\frac{p}{1-p}=w^Tx$ (log-odds linear in features).

% =========================
\section*{Logistic: Hessian \& Convexity}

Log-loss: $\ell(\beta)=\sum_i[\log(1+e^{z_i})-y_iz_i]$, $z_i=x_i^T\beta$.

Gradient: $\nabla\ell=X^T(p-y)$, $p_i=\sigma(z_i)$.

Hessian: $\nabla^2\ell=X^TWX$, $W=\text{diag}(p_i(1-p_i))$.

$p_i(1-p_i)\ge0$ $\Rightarrow$ $\nabla^2\ell\succeq0$ $\Rightarrow$ convex.

Boundary: $\sigma(\beta^Tx)=0.5\iff\beta^Tx=0$ (linear).

Logit: $\log\frac{p}{1-p}=x^T\beta$.


% =========================
\section*{Bias-Variance}
$\mathbb E[(y-\hat f)^2]=\underbrace{(\mathbb E[\hat f]-f)^2}_{\text{Bias}^2}+\underbrace{\mathbb E[(\hat f-\mathbb E[\hat f])^2]}_{\text{Var}}+\sigma^2$.
Bias: error from wrong assumptions (underfitting). Var: sensitivity to training data (overfitting). $\sigma^2$: irreducible noise. High $\lambda$ $\Rightarrow$ high bias, low var. Low $\lambda$ $\Rightarrow$ low bias, high var.

% =========================
\section*{SST/SSM/SSE (Complete)}

$SST=\sum(y_i-\bar y)^2$, $SSE=\sum(y_i-\hat y_i)^2$, $SSM=\sum(\hat y_i-\bar y)^2$.

$SST=SSM+SSE$.

$R^2=\frac{SSM}{SST}=1-\frac{SSE}{SST}$.

$R^2=1$: perfect fit. $R^2=0$: no better than mean.

Correlation: $r=\hat\beta_1\frac{\sigma_x}{\sigma_y}$, $R^2=r^2$ (simple).


% =========================
\section*{Confusion Matrix + Metrics}

Binary classification:

$\begin{array}{c|cc}
 & \text{Pred +} & \text{Pred -} \\
\hline
\text{Actual +} & TP & FN \\
\text{Actual -} & FP & TN
\end{array}$

Accuracy: $\frac{TP+TN}{TP+TN+FP+FN}$.

Precision: $\frac{TP}{TP+FP}$.

Recall: $\frac{TP}{TP+FN}$.

$F1=\frac{2PR}{P+R}=\frac{2TP}{2TP+FP+FN}$.

\textbf{$R^2$:}
$SST=\sum (y_i-\bar y)^2$,
$SSE=\sum (y_i-\hat y_i)^2$.

$R^2=1-\frac{SSE}{SST}=\frac{SST-SSE}{SST}$.

% =========================
\section*{Ridge (Complete)}

$\hat\beta=(X^TX+\lambda I)^{-1}X^Ty$.

$\lambda>0$: $X^TX+\lambda I\succ0$ always invertible.

Handles collinearity; shrinks coeff.; not sparse.

MAP: Gaussian prior $\beta\sim\mathcal N(0,\tau^2 I)\Rightarrow$ Ridge.

$\lambda\to0$: converges to min-norm OLS $X^+y$.


% =========================
\section*{LASSO (Complete)}

$\hat\beta=\arg\min\|y-X\beta\|^2+\lambda\|\beta\|_1$.

Sparse (feature selection); no closed form.

MAP: Laplace prior $\beta_j\sim\text{Lap}(0,1/\lambda)\Rightarrow$ LASSO.

Ridge vs LASSO: Ridge shrinks all, LASSO zeros some.


% =========================
\section*{Collinearity}

If $x_{\cdot,2}=cx_{\cdot,1}$: $X^TX$ singular.

$\forall v\in\text{null}(X)$: $J(\beta+v)=J(\beta)$.

$\Rightarrow$ infinitely many OLS solutions.

Ridge adds $\lambda I$: unique solution for any $\lambda>0$.


% =========================
\section*{Robust Regression}

$r_i=y_i-x_i^T\beta$.

L2: $\rho=r^2$. Gaussian MLE. Sensitive to outliers.

L1/LAD: $\rho=|r|$. Laplace MLE. Robust.

Huber ($k>0$):
$\rho_k(r)=\begin{cases}\frac12r^2,&|r|\le k,\\k(|r|-\frac{k}{2}),&|r|>k.\end{cases}$

Quadratic near 0, linear in tails.

Gaussian noise $\Rightarrow$ L2. Laplace noise $\Rightarrow$ L1.


% =========================
\section*{Cross-Validation}

Holdout: 80/20 split. Fast, high-variance estimate.

$k$-fold CV: $k$ folds; each fold = test once.
$\text{CV}=\frac1k\sum_{j=1}^k\text{err}_j$. Default $k=5,10$.

LOOCV ($k=n$): $\text{CV}=\frac1n\sum\!\left(\frac{e_i}{1-h_{ii}}\right)^2$.

Select hyperparameter by CV error, NOT train error.


% =========================
\section*{Hat Matrix Properties}

$H=X(X^TX)^{-1}X^T$.

$\text{tr}(H)=\text{tr}((X^TX)^{-1}X^TX)=\text{tr}(I_p)=p$.

$0\le h_{ii}\le1$.

Large $h_{ii}$: point $i$ strongly controls its own fit (high leverage).

LOOCV: $e_i^{(-i)}=\frac{e_i}{1-h_{ii}}$.


% =========================
\section*{Feature Engineering}

Polynomial: $\phi(x)=(1,x,x^2,\dots,x^p)$.
$\Rightarrow$ linear in $\phi$-space, nonlinear in $x$.

Standardization: $\tilde x=\frac{x-\bar x}{\sigma}$.
Equalizes scale; improves SGD/perceptron convergence.

XOR: not separable in $\mathbb R^2$; separable with $(x_1,x_2,x_1x_2)$.


% =========================
\section*{Cost-Sensitive Threshold}

$\eta(x)=P(Y=1|x)$.

Cost(predict 1): $C_{FP}(1-\eta)$.
Cost(predict 0): $C_{FN}\eta$.

Predict $1$ iff $\eta\ge\dfrac{C_{FP}}{C_{FP}+C_{FN}}$.

Default 0.5 $\iff$ $C_{FP}=C_{FN}$.

High $C_{FP}$ (expensive FP): raise threshold.
High $C_{FN}$ (expensive FN): lower threshold.


% =========================
% GROUP 2: PROBABILISTIC
% =========================
\section*{Bayes Classifier}

\textbf{Bayes Theorem:}
$P(Y|X)=\frac{P(X|Y)P(Y)}{P(X)}$.

\textbf{Optimal Bayes Rule:}

$h^*(x)=\arg\max_{y\in\mathcal Y} P(Y=y|X=x)$.

Binary case ($Y=\{0,1\}$):

$\eta(x)=P(Y=1|x)$,

$h^*(x)=\mathbf{1}[\eta(x)>0.5]$.

\textbf{Why Optimal?}

Conditional risk at $x$:
$R(h|x)=1-P(h(x)|x)$.

Choosing class with largest posterior
minimizes $R(h|x)$ for every $x$.

Thus for any classifier $g$:
$R(h^*)\le R(g)$.

\textbf{Bayes Risk:}

Minimum achievable error:

$R^*=\mathbb E_X[1-\max_k P(Y=k|X)]$.

Represents irreducible error due to class overlap.

\textbf{Key Intuition:}
• Uses full distribution knowledge.
• Always chooses most probable label.
• No classifier can do better pointwise.


% =========================
\section*{Naive Bayes}

\textbf{Definition:}
Generative classifier using independence assumption.

Assumption:
$X_1,\dots,X_n$ conditionally independent given $Y$.

Posterior:

$P(Y|X_1,\dots,X_n)
\propto
P(Y)\prod_{i=1}^n P(X_i|Y)$.

Prediction:

$h(x)=\arg\max_y P(Y=y)\prod_i P(x_i|Y=y)$.

\textbf{Why "Naive"?} Assumes features independent given class (often violated, yet works well).

\textbf{Gaussian NB}: $P(x_j|y)=\mathcal N(\mu_{jy},\sigma_{jy}^2)$; estimate $\mu_{jy},\sigma_{jy}^2$ from class samples.

\textbf{Categorical NB}: $P(x_j=v|y)=\frac{\#(x_j=v,y)+\alpha}{\#y+\alpha|\mathcal V|}$ (Laplace smoothing $\alpha>0$ avoids zero probs).

Use log-space: $\log P(y)+\sum_j\log P(x_j|y)$ to avoid underflow.

\textbf{Pros}: $O(nd)$ training, handles high-$d$, missing features, small data.


% =========================
\section*{Generative vs Discriminative}

\textbf{Generative Methods:}
Model joint distribution $P(X,Y)$.
Can generate data.
Example: Naive Bayes.

\textbf{Discriminative Methods:}
Model conditional $P(Y|X)$
or decision boundary directly.
Example: SVM.

\textbf{Key Difference:}

Generative:
learn data distribution.

Discriminative:
learn boundary between classes.

Generative can simulate new samples.
Discriminative often achieve lower classification error.


% =========================
\section*{Bayesian Linear Regression}

Model: $y=Xw+\epsilon$, $\epsilon\sim\mathcal N(0,\sigma^2 I)$.

Likelihood: $p(y|X,w)=\mathcal N(Xw,\sigma^2 I)$.

Prior: $p(w)=\mathcal N(m_0,S_0)$.

\textbf{Posterior} (Gaussian conjugacy):
$S_N=\!\left(S_0^{-1}+\sigma^{-2}X^TX\right)^{-1}$,\quad $m_N=S_N\!\left(S_0^{-1}m_0+\sigma^{-2}X^Ty\right)$.

\textbf{Special case} ($m_0=0$, $S_0=\alpha^{-1}I$):
$m_N=(X^TX+\alpha\sigma^2 I)^{-1}X^Ty$ $\Rightarrow$ \textbf{Ridge}!

\textbf{Posterior predictive} (at $x^*$):
$p(y^*|x^*,X,y)=\mathcal N(\mu^*,(\sigma^*)^2)$,
$\mu^*=x^{*T}m_N$, $(\sigma^*)^2=\sigma^2+x^{*T}S_N x^*$.

$(\sigma^*)^2$: aleatoric ($\sigma^2$) + epistemic ($x^{*T}S_Nx^*$).

Large $\alpha$ (strong prior, small variance) $\Rightarrow$ more regularization.

MAP estimate $= m_N$ (matches Ridge/LASSO depending on prior).


% =========================
% GROUP 3: KNN
% =========================
\section*{K-Nearest Neighbours (KNN)}

\textbf{Motivation:}
• Works well when data is large and complex to model.
• Simple, intuitive, non-parametric.
• No assumption on underlying distribution.
• Works for classification and regression.

\textbf{Algorithm:}
Given training set $S=\{(x_i,y_i)\}_{i=1}^m$.

For new input $x$:
1) Compute distance $d(x,x_i)$ for all $i$.
2) Select $K$ nearest neighbours.
3) Classification: majority vote.
4) Regression: average of neighbours.

Prediction:
$h(x)=\arg\max_y \sum_{i\in N_K(x)} \mathbf{1}[y_i=y]$.

\textbf{Distance Measures:}

Manhattan:
$d(x,y)=\|x-y\|_1=\sum_{i=1}^n |x_i-y_i|$.

Euclidean:
$d(x,y)=\|x-y\|_2=\sqrt{\sum (x_i-y_i)^2}$.

Minkowski:
$d(x,y)=\|x-y\|_p=(\sum |x_i-y_i|^p)^{1/p}$.

$p=1 \Rightarrow$ Manhattan,
$p=2 \Rightarrow$ Euclidean.

\textbf{Properties:}
• Lazy learner (no explicit training phase).
• Decision boundary can be highly non-linear.
• As $n\to\infty$, 1-NN error $\le 2R^*$.

\textbf{Limitations:}
• Curse of dimensionality.
• Slow for large datasets (distance to all points).
• Sensitive to noise.
• Memory intensive.


% =========================
\section*{Feature Scaling in KNN}

Without scaling, large-magnitude features dominate Euclidean distance.

Min-Max: $x'_j=\frac{x_j-x_{j,\min}}{x_{j,\max}-x_{j,\min}}\in[0,1]$.

Z-score: $x'_j=\frac{x_j-\mu_j}{\sigma_j}$ (zero mean, unit var).

Scaling equalizes feature influence $\Rightarrow$ NN assignments change.
Must scale \textbf{before} computing distances; fit scaler on train set only.


% =========================
\section*{Bayes vs KNN}

Bayes:
• Optimal if true distribution known.
• Requires estimating $P(Y|X)$.

KNN:
• Distribution-free.
• Approximates Bayes decision boundary
as sample size grows.

1-NN error bound:
$R_{1NN}\le 2R^*$.

% =========================
\section*{1NN vs Bayes}

Asked: show $R_{1NN}\le 2R^*$ as $n\to\infty$.

Let $\eta(x)=P(Y=1|x)$.

Bayes classifier predicts $1$ if $\eta\ge 1/2$.

Bayes error at $x$:
$R^*(x)=\min(\eta,1-\eta)$.

As $n\to\infty$, nearest neighbor $X_{(1)}\to x$.
Hence $Y_{(1)}\sim Bern(\eta(x))$ and is independent of $Y$ given $x$.

1NN predicts $Y_{(1)}$.

\[
R_{1NN}(x)=P(Y\neq Y_{(1)}|x)
\]

\[
=\eta(1-\eta)+(1-\eta)\eta
=2\eta(1-\eta)
\]

If $\eta\le1/2$, then $R^*(x)=\eta$.

\[
R_{1NN}(x)=2R^*(1-R^*)\le2R^*(x)
\]

Taking expectation over $x$:

\[
R_{1NN}\le 2R^*
\]


% =========================
\section*{Curse of Dimensionality}

$X,Y\sim\text{Unif}([0,1]^d)$. By LLN:

$\frac{\|X-Y\|^2}{d}\to\frac16$.

Distances concentrate around $\sqrt{d/6}$.

$\frac{D_{\max}-D_{\min}}{D_{\min}}\to0$ as $d\to\infty$.

NN becomes unreliable (all points equidistant).

Need $\sim N^d$ samples to cover $[0,1]^d$.


% =========================
% GROUP 4: LINEAR CLASSIFIERS
% =========================
\section*{Linear Separability}

Training set: $S=\{(x_i,y_i)\}_{i=1}^m$, $x_i\in\mathbb{R}^d$, $y_i\in\{-1,+1\}$.

Data is linearly separable if $\exists (w,b)$ such that:

$y_i(\langle w,x_i\rangle+b) > 0 \quad \forall i$.

Decision rule:
$\hat y=\text{sign}(\langle w,x\rangle+b)$.

% =========================
\section*{Margin}

Margin of hyperplane $(w,b)$:

$\gamma = \min_i \frac{| \langle w,x_i\rangle + b |}{\|w\|}$.

If $\|w\|=1$:

Distance of $x$ to hyperplane = $|\langle w,x\rangle+b|$.

Geometric margin:
$\gamma = \min_i y_i(\langle w,x_i\rangle+b)$.

Larger margin implies better generalization.

% =========================
\section*{Multiclass Perceptron}

Assume $\exists w^*$ s.t.
$\langle w^*,\Psi(x_i,y_i)\rangle \ge \langle w^*,\Psi(x_i,y)\rangle+1$.

Update:
$w^{t+1}=w^t+\Psi(x_i,y_i)-\Psi(x_i,y)$.

Progress:
$\langle w^*,w^{t+1}\rangle
=\langle w^*,w^t\rangle
+\langle w^*,\Psi(x_i,y_i)-\Psi(x_i,y)\rangle
\ge \langle w^*,w^t\rangle+1$.

After $T$: $\langle w^*,w^T\rangle\ge T$.

Norm growth:
$||w^{t+1}||^2
\le ||w^t||^2+R^2$.

Thus $||w^T||^2\le TR^2$.

Combine:
$T \le ||w^*||\,||w^T||
\le ||w^*||\sqrt{TR^2}$.

$T\le R^2||w^*||^2$.

% =========================
\section*{Perceptron as Subgradient}

Let $S=\{(x_i,y_i)\}_{i=1}^m$, $y_i\in\{\pm1\}$.
Assume $\exists w$ s.t. $y_i\langle w,x_i\rangle \ge 1$ for all $i$.

Define
$f(w)=\max_{i\in[m]}(1-y_i\langle w,x_i\rangle)$.

Then $f(w)\ge0$ and $f(w)=0$ iff $y_i\langle w,x_i\rangle\ge1$ $\forall i$.
Thus $\min_{\|w\|\le\|w^*\|} f(w)=0$.

If $f(w)<1$, then
$1-y_i\langle w,x_i\rangle<1$
$\Rightarrow y_i\langle w,x_i\rangle>0$,
so $w$ separates the data.

Subgradient:
Let $i^*=\arg\max_i (1-y_i\langle w,x_i\rangle)$.
If $f(w)>0$ (constraint violated),

$\partial f(w)= -y_{i^*}x_{i^*}$

since $f$ is max of affine functions.

Subgradient descent with step $\eta$:

$w^{t+1}=w^t-\eta(-y_{i^*}x_{i^*})
= w^t+\eta y_{i^*}x_{i^*}$.

This is exactly the Perceptron update
(applied when $y_i\langle w,x_i\rangle \le 0$).

Thus Perceptron = subgradient descent on $f$.


% =========================
\section*{Margin Bound (Perceptron)}

Assume:
Data linearly separable with margin $\gamma$.
$\exists w^*$ s.t. $y_i\langle w^*,x_i\rangle \ge \gamma$.
Also $||x_i||\le\rho$.

Perceptron update:
$w^{t+1}=w^t+y_ix_i$ (on mistake).

1) Progress toward $w^*$:

$\langle w^*,w^{t+1}\rangle
=\langle w^*,w^t\rangle
+y_i\langle w^*,x_i\rangle
\ge
\langle w^*,w^t\rangle+\gamma$.

After $T$ mistakes:

$\langle w^*,w^T\rangle \ge T\gamma$.

2) Norm growth:

$||w^{t+1}||^2
=||w^t||^2+||x_i||^2
+2y_i\langle w^t,x_i\rangle$.

Since mistake:
$y_i\langle w^t,x_i\rangle \le 0$.

Thus

$||w^{t+1}||^2 \le ||w^t||^2+\rho^2$.

After $T$ updates:

$||w^T||^2 \le T\rho^2$.

Now by Cauchy--Schwarz:

$T\gamma
\le
\langle w^*,w^T\rangle
\le
||w^*||\,||w^T||
\le
||w^*||\sqrt{T\rho^2}$.

Divide by $\sqrt{T}$:

$\sqrt{T}\gamma \le ||w^*||\rho$.

Thus

$T \le (\rho/\gamma)^2$.

Finite mistake bound proved.


% =========================
\section*{Hard-SVM Rule}

Goal: find separator with largest margin.

\[
\max_{w,b} \min_i y_i(\langle w,x_i\rangle+b)
\]

Equivalent constrained form:

\[
\min_{w,b} \frac12\|w\|^2
\quad \text{s.t. } y_i(\langle w,x_i\rangle+b)\ge1
\]

Maximizing margin iff minimizing $\|w\|$.

% =========================
\section*{Hard-SVM Lemma}

Let $(w^*,b^*)$ achieve margin $\gamma^*$.

Then scaling:

$\left(\frac{w^*}{\gamma^*},\frac{b^*}{\gamma^*}\right)$

satisfies constraints $y_i(\langle w,x_i\rangle+b)\ge1$.

Thus:

$\|w_0\|\le \frac1{\gamma^*}$.

Final normalized solution:

$\hat w=\frac{w_0}{\|w_0\|}$.

Hence solves original margin maximization.

% =========================
\section*{Homogeneous Case}

If $b=0$:

\[
\min_w \frac12\|w\|^2
\quad \text{s.t. } y_i\langle w,x_i\rangle\ge1
\]

Separating hyperplanes pass through origin.

% =========================
\section*{Hard-Margin SVM}

Asked: find $(w,b)$, support vectors, effect of removing $(5,5)$.

Data separable by line $x_1+x_2=5$.

Let $w=(a,a)$, $b=-5a$.

Margin condition:
For support vectors,
$y_i(w^Tx_i+b)=1$.

For $(4,2)$:
$a(6)-5a=a=1 \Rightarrow a=1/2$.

Hence:
$w=(\tfrac12,\tfrac12)$,
$b=-\tfrac52$.

Support vectors satisfy equality:
$(4,2),(2,4),(1,2),(2,1)$.

For $(5,5)$:
$5+5-5=5>1$,
strict inequality $\to$ not SV.

Removing it does not change solution.


% =========================
\section*{Fritz John Optimality Insight}

If $w^*$ minimizes $f(w)$ under constraints $g_i(w)\le0$:

$\nabla f(w^*) + \sum_{i\in I}\alpha_i\nabla g_i(w^*)=0$,

$I=\{i: g_i(w^*)=0\}$.

Applying to Hard-SVM yields:

$w=\sum_i \alpha_i y_i x_i$.

% =========================
\section*{Dual of Hard-SVM}

From optimality condition:

$w=\sum_i \alpha_i y_i x_i$.

Substitute into primal:

\[
\max_{\alpha\ge0}
\sum_i \alpha_i
-
\frac12\sum_{i,j}
\alpha_i\alpha_j y_i y_j
\langle x_i,x_j\rangle
\]

Dual depends only on inner products.

Enables kernel methods.

% =========================
\section*{Kernel Trick (SVM Dual)}

Dual SVM depends only on inner products:

\[
\sum_{i,j}\alpha_i\alpha_j y_i y_j
\langle x_i,x_j\rangle
\]

Replace:

\[
\langle x_i,x_j\rangle
\rightarrow
K(x_i,x_j)
\]

Common kernels:

Linear:
$K(x,z)=\langle x,z\rangle$

Polynomial:
$(\langle x,z\rangle+c)^p$

RBF:
$e^{-\gamma\|x-z\|^2}$

Enables nonlinear decision boundaries.


% =========================
\section*{Soft-SVM (Slack Variables)}

For non-separable data introduce $\xi_i\ge0$:

\[
y_i(\langle w,x_i\rangle+b)\ge1-\xi_i
\]

Optimization:

\[
\min_{w,b,\xi}
\lambda\|w\|^2+\frac1m\sum_i \xi_i
\]

Trade-off controlled by $\lambda$.

Large $\lambda$ $\to$ behaves like hard SVM.

% =========================
\section*{Hinge Loss Formulation}

Hinge loss:

$\ell_{hinge}((w,b),(x,y))=\max(0,1-y(\langle w,x\rangle+b))$.

Soft-SVM:

\[
\min_{w,b}
\lambda\|w\|^2
+
\frac1m\sum_i
\max(0,1-y_i f(x_i))
\]

Convex optimization problem.

% =========================
\section*{Support Vectors}

At optimum (Hard SVM):

$w_0 = \sum_{i\in I} \alpha_i x_i$

where $I=\{i: y_i(\langle w_0,x_i\rangle+b_0)=1\}$.

Only points exactly at margin contribute.

These are called support vectors.

They lie in linear span of training examples.

% =========================
\section*{Hard vs Soft SVM}

Hard-SVM:

$\min_{w,b} ||w||^2$
s.t. $y_i(w^Tx_i+b)\ge1$.

Soft-SVM:

$\min_{w,b,\xi}
||w||^2+\lambda\sum_i\xi_i$
s.t. $y_i(w^Tx_i+b)\ge1-\xi_i$,
$\xi_i\ge0$.

Claim:
$\exists\lambda>0$ s.t. for all separable datasets,
hard and soft return same $w$.

Refutation:

Even if data separable,
finite $\lambda$ trades margin size vs slack penalty.

For different datasets,
required $\lambda$ to force $\xi_i=0$ differs.

Example:
Two separable datasets with very different margins.
Same fixed $\lambda$ cannot simultaneously enforce zero slack
for both while preserving identical solution.

Only in limit $\lambda\to\infty$
does soft-SVM reduce to hard-SVM.

Hence no universal $\lambda$ works.
Claim false.

% =========================
\section*{Soft-Margin SVM (1D)}

Asked: number of SV when $C=0$ and $C\to\infty$.

Data separable.

If $C\to\infty$:
No slack allowed $\to$ hard margin solution.

Closest opposite points: $-1$ and $1$.

Decision boundary at midpoint $0$.

$w=1$, $b=0$.

Support vectors: $-1$ and $1$ (2 SV).

If $C=0$:
Objective reduces to minimizing $\frac12||w||^2$.

Minimum at $w=0$.

All constraints violated $\to$ all 5 points are support vectors.


% =========================
\section*{Why $\alpha_i=C$}

Asked: show $\alpha_i=C \implies$ inside margin or misclassified.

KKT conditions:

$\alpha_i(1-y_if(x_i)-\xi_i)=0$,

$\beta_i\xi_i=0$,

$\alpha_i+\beta_i=C$.

If $\alpha_i=C$,
then $\beta_i=0$.

From complementarity:

$y_if(x_i)=1-\xi_i$.

If $\xi_i=0$:
$y_if=1$ (on margin).

If $\xi_i>0$:
$y_if<1$ (inside margin).

If $\xi_i>1$:
$y_if<0$ (misclassified).

Hence $\alpha_i=C$ implies on/inside margin or misclassified.


% =========================
\section*{Convexity + Non-Differentiability}

Asked:
(a) prove convexity of $\ell(w)=|y-\max(0,w^Tx)|$,
(b) find non-diff points.

Let $f(w)=\max(0,w^Tx)$.

Since $w^Tx$ is affine and max of affine functions is convex,
$f(w)$ is convex.

Define $g(w)=y-f(w)$.

Affine minus convex $\to$ convex.

Since $|z|=\max(z,-z)$ (max of convex functions),
$\ell(w)$ is convex.

Non-differentiability:

ReLU kink at $w^Tx=0$.

Absolute value kink when $y-f(w)=0$.

Thus non-diff when:
$w^Tx=0$ or $y=\max(0,w^Tx)$.


% =========================
\section*{Weak Duality}

Claim:
$\min_{x\in\mathcal X}\max_{y\in\mathcal Y} f(x,y)
\ge
\max_{y\in\mathcal Y}\min_{x\in\mathcal X} f(x,y)$.

Proof:

For any fixed $(x,y)$:

$\min_{x'} f(x',y) \le f(x,y) \le \max_{y'} f(x,y')$.

Taking $\max_y$ on the left inequality:

$\max_y \min_x f(x,y) \le \max_y f(x,y)$.

Now take $\min_x$ on RHS:

$\max_y \min_x f(x,y)
\le
\min_x \max_y f(x,y)$.

Hence proved.

No convexity assumptions required.


% =========================
\section*{Overfitting vs Regularization}

High bias $\to$ underfitting.
High variance $\to$ overfitting.

Regularization effect:

Increase $\lambda$ $\to$
• increases bias
• reduces variance
• smoother model

Trade-off controlled by $\lambda$.


% =========================
\section*{VC Dimension}

$\text{VC}(\mathcal H)=$ max $m$ s.t. $\exists m$ points shattered.

Linear classifiers in $\mathbb R^d$: $\text{VC}=d+1$.

Decision trees depth $k$: $O(2^k\log d)$; unrestricted: $\infty$.

SVM RBF (no reg.): $\text{VC}=\infty$.

Generalization bound:
$R(h)\le\hat R(h)+O\!\left(\!\sqrt{\frac{\text{VC}\log m}{m}}\right)$.


% =========================
\section*{Perceptron vs Logistic vs SVM}

Perceptron: no convex obj.; converges only if separable; no margin guarantee; sensitive to outliers.

Logistic: convex log-loss; probabilistic output; global optimum; no explicit margin.

Hard-SVM: max margin; requires separability; unique; sensitive to noise.

Soft-SVM: hinge loss; handles non-separable; controlled by $\lambda$.


% =========================
% GROUP 5: MULTICLASS
% =========================
\section*{One-vs-All (OvA)}

Given $S=\{(x_t,y_t)\}_{t=1}^m$, label set $\mathcal Y=\{1,\dots,k\}$.

For each class $i$:
Construct binary dataset
$S_i=\{(x_t, +1)$ if $y_t=i$, $(x_t,-1)$ otherwise$\}$.

Train binary classifier $h_i$.

Final predictor:
$h(x)\in \arg\max_{i\in\mathcal Y} h_i(x)$.

Issue:
Each classifier sees highly imbalanced data.
May output trivial all-negative classifier for rare class.


% =========================
\section*{All-Pairs (One-vs-One)}

For each pair $i<j$:

Construct
$S_{i,j}=\{(x_t,+1)$ if $y_t=i$,
$(x_t,-1)$ if $y_t=j\}$.

Train $h_{i,j}$.

Prediction by voting:

$h(x)\in\arg\max_i \sum_{j\neq i} h_{i,j}(x)$.

Requires $k(k-1)/2$ classifiers.

Still may suffer inconsistency across pairwise decisions.


% =========================
\section*{Failure of Reduction}

Example: 3 classes in $\mathbb R^2$ in disjoint balls.
Probabilities: 40\%, 20\%, 40\%.

OvA classifier for class 2 vs rest
may choose all-negative solution
$\to$ misclassifies all class-2 points.

But joint linear predictor can perfectly separate.

Conclusion:
Reduction methods can fail
even if perfect linear separator exists.


% =========================
\section*{Linear Multiclass Predictor}

Single parameter vector:

\[
h(x)=\arg\max_{y\in\{1,\dots,k\}}
\langle w,\Psi(x,y)\rangle
\]

$\Psi(x,y)$: class-sensitive feature map.

Multi-vector construction:

$\Psi(x,y)$ = concat of $k$ copies of $x$,
only $y$-block nonzero.

Then:
\[
\langle w,\Psi(x,y)\rangle=\langle w_y,x\rangle
\]

Each class has weight vector $w_y$.


% =========================
\section*{Multiclass Hinge Loss}

For example $(x,y)$:

\[
\ell(w;x,y)
=
\max_{y'\neq y}
\Big[
\Delta(y',y)
+
\langle w,
\Psi(x,y')-\Psi(x,y)
\rangle
\Big]
\]

$\Delta$ = cost (0--1 or cost-sensitive).

Convex surrogate of 0--1 loss.


% =========================
\section*{Multiclass SVM}

\[
\min_w
\lambda||w||^2
+
\frac1m
\sum_{i=1}^m
\ell(w;x_i,y_i)
\]

Convex optimization problem.

Handles cost-sensitive errors via $\Delta$.


% =========================
\section*{Multiclass SGD}

Let

\[
y^*=
\arg\max_{y'}
\Big[
\Delta(y',y)
+
\langle w,
\Psi(x,y')-\Psi(x,y)
\rangle
\Big]
\]

Subgradient:
\[
g=\Psi(x,y^*)-\Psi(x,y)
\]

Update:
\[
w_{t+1}=w_t-\eta_t g
\]

Scales to large datasets.
Has generalization guarantees.


% =========================
\section*{Advantages of Joint Multiclass}

• Single model captures class relations
• Avoids inconsistencies of reductions
• Supports cost-sensitive learning

If realizable:
use linear programming or perceptron.

If non-realizable:
use hinge loss + convex optimization.


% =========================
\section*{Ranking}

Goal: order items by relevance (search, fraud detection).

Linear model: score $s_j=\langle w,x_j\rangle$; sort by score descending.

\textbf{Losses}: Kendall-Tau = pairwise disagreements (0--1, NP-hard); use surrogate pairwise hinge. NDCG = normalized discounted cumulative gain (emphasizes top positions). Learn $w$ via surrogate loss + SGD.


% =========================
% GROUP 6: OPTIMIZATION
% =========================
\section*{Gradient Descent (GD)}

Goal: minimize risk $L_D(w)=\mathbb E_{z\sim D}[\ell(w,z)]$.

Gradient:
$\nabla f(w)=\left(\frac{\partial f}{\partial w_1},\dots,\frac{\partial f}{\partial w_d}\right)$.

Update rule:
$w^{(t+1)}=w^{(t)}-\eta\nabla f(w^{(t)})$.

$\eta>0$ is learning rate.

\textbf{Taylor view:}
$f(u)\approx f(w)+\langle u-w,\nabla f(w)\rangle$.

For convex $f$:
$f(u)\ge f(w)+\langle u-w,\nabla f(w)\rangle$.

Equivalent proximal view:

$w^{(t+1)}=
\arg\min_w
\left[\frac12\|w-w^{(t)}\|^2
+\eta\langle w-w^{(t)},\nabla f(w^{(t)})\rangle\right]$.


% =========================
\section*{Subgradients}

GD requires differentiability.
Extend to convex non-differentiable functions.

Definition:
$v\in\partial f(w)$ if
$\forall u,\;
f(u)\ge f(w)+\langle u-w,v\rangle$.

$\partial f(w)$ = subdifferential set.

If $f$ differentiable:
$\partial f(w)=\{\nabla f(w)\}$.

Example:
$f(x)=|x|$.

$\partial f(x)=
\begin{cases}
\{1\}, & x>0 \\
\{-1\}, & x<0 \\
[-1,1], & x=0
\end{cases}$.

\textbf{Subgradient of max:}

If $g(w)=\max_i g_i(w)$ and $j\in\arg\max_i g_i(w)$,
then $\nabla g_j(w)\in\partial g(w)$.

\textbf{Hinge loss:}

$\ell(w)=\max(0,1-y\langle w,x\rangle)$.

Subgradient:
$v=
\begin{cases}
0, & 1-y\langle w,x\rangle\le0 \\
-yx, & 1-y\langle w,x\rangle>0
\end{cases}$.


% =========================
\section*{Lipschitz Convex Functions}

$f$ is $\rho$-Lipschitz iff
$\|v\|\le\rho$ for all $v\in\partial f(w)$.

Controls gradient magnitude.


% =========================
\section*{Stochastic Gradient Descent (SGD)}

We cannot compute $\nabla L_D(w)$ (unknown $D$).

SGD uses random direction $v_t$ s.t.

$\mathbb E[v_t|w^{(t)}]\in\partial f(w^{(t)})$.

Update:
$w^{(t+1)}=w^{(t)}-\eta v_t$.

Output average:
$\bar w=\frac1T\sum_{t=1}^T w^{(t)}$.

Key idea:
Unbiased estimate of gradient.

GD: exact gradient.
SGD: noisy but cheaper.


% =========================
\section*{SGD Algorithm}

Init $w^{(1)}=0$. For $t=1,\dots,T$: pick random $v_t$ with $\mathbb E[v_t|w^{(t)}]\in\partial f(w^{(t)})$; update $w^{(t+1)}=w^{(t)}-\eta_t v_t$. Return $\bar w=\frac1T\sum_t w^{(t)}$.

Variable step: $\eta_t=\frac{B}{\rho\sqrt{t}}$ (decaying) or fixed $\eta=\sqrt{B^2/(\rho^2 T)}$.

\textbf{Convergence}: $\mathbb E[f(\bar w)]-f(w^*)\le\frac{B\rho}{\sqrt{T}}$. Need $T\ge B^2\rho^2/\epsilon^2$ for $\epsilon$-suboptimality.

% =========================
\section*{Strongly Convex SGD}

$f$ is $\lambda$-strongly convex: $f(u)\ge f(w)+\langle v,u-w\rangle+\frac\lambda2\|u-w\|^2$.

Use $\eta_t=\frac1{\lambda t}$. Convergence: $O(\rho^2/(\lambda T))$ (faster than general case).

Projection step: $w^{(t+1)}=\Pi_{\mathcal H}(w^{(t)}-\eta_t v_t)$ where $\Pi_{\mathcal H}(w)=\arg\min_{u\in\mathcal H}\|u-w\|$.


% =========================
\section*{Projected SGD}

Feasible set:
$\mathcal H=\{w:\|w\|\le B\}$.

Step 1:
$w^{(t+\frac12)}=w^{(t)}-\eta v_t$.

Step 2 (projection):
$w^{(t+1)}=\arg\min_{w\in\mathcal H}
\|w-w^{(t+\frac12)}\|$.

Projection lemma:
For projection $v$ of $w$ onto convex $\mathcal H$:

$\|v-u\|^2\le\|w-u\|^2$
for all $u\in\mathcal H$.

Ensures distance to optimum does not increase.


% =========================
\section*{Convex-Lipschitz-Bounded Case}

Assume:
• $\|w\|\le B$
• $\|v\|\le\rho$

If

$T\ge \frac{B^2\rho^2}{\epsilon^2}$

and

$\eta=\sqrt{\frac{B^2}{\rho^2 T}}$,

then

$\mathbb E[L_D(\bar w)]
\le
\min_{w\in\mathcal H}L_D(w)+\epsilon$.

Sample complexity matches
regularized ERM order.


% =========================
\section*{GD vs SGD}

GD:
• Uses full gradient.
• Expensive for large datasets.

SGD:
• Uses single random example.
• Faster per iteration.
• Converges in expectation.

SGD average iterate stabilizes noise.

% =========================
\section*{Empirical Risk Minimization (ERM)}

True risk:
\[
L_D(w)=\mathbb E_{(x,y)\sim D}[\ell(w,(x,y))]
\]

Empirical risk:
\[
L_S(w)=\frac1m\sum_{i=1}^m \ell(w,(x_i,y_i))
\]

ERM:
\[
\hat w=\arg\min_w L_S(w)
\]

Regularized ERM:
\[
\hat w=\arg\min_w L_S(w)+\lambda R(w)
\]

Used in SVM, Logistic, Ridge.


% =========================
\section*{0--1 vs Surrogate Loss}

0--1 loss:
\[
\ell_{0-1}(w,(x,y))
=
\mathbf 1\{y f(x)\le0\}
\]

Non-convex $\to$ hard to optimize.

Use convex upper bounds:

Hinge:
\[
\max(0,1-yf(x))
\]

Logistic:
\[
\log(1+e^{-y f(x)})
\]

Surrogates enable convex optimization.


% =========================
\section*{Logistic Margin Loss}

Binary logistic model:

\[
P(y=1|x)=\sigma(\langle w,x\rangle)
\]

Loss:
\[
\ell(w,(x,y))
=
\log(1+e^{-y\langle w,x\rangle})
\]

Convex and smooth surrogate of 0--1.


% =========================
\section*{Sherman-Morrison + LOOCV}

Asked:
1) prove inverse identity,
2) derive LOOCV formula.

Sherman-Morrison:

\[
(A-uu^T)^{-1}
=
A^{-1}
+
\frac{A^{-1}uu^TA^{-1}}
{1-u^TA^{-1}u}
\]

Verify by multiplying RHS with $(A-uu^T)$ to obtain identity.

OLS:

$\hat\beta=(X^TX)^{-1}X^Ty$.

Hat matrix:
$H=X(X^TX)^{-1}X^T$.

Predictions:
$\hat y=Hy$.

Residual:
$e_i=y_i-\hat y_i$.

Leaving out $i$:
Apply Sherman-Morrison to
$X^TX-x_ix_i^T$.

Result:

$\hat y_i^{(-i)}
=
\hat y_i
-
\frac{h_i e_i}{1-h_i}$.

Prediction error:

$y_i-\hat y_i^{(-i)}
=
\frac{e_i}{1-h_i}$.

LOOCV:

$CV=\frac1n\sum
\left(\frac{y_i-\hat y_i}{1-h_i}\right)^2$.

Single model fit sufficient.

% =========================
% GROUP 7: ENSEMBLE METHODS
% =========================
\section*{Decision Trees}

ID3 (greedy top-down): start at root; split on feature maximizing IG; recurse; stop if pure/max-depth/min-samples; leaf = majority class.

\textbf{Impurity} (binary $p$ = fraction positive):
Entropy: $H(p)=-p\log p-(1-p)\log(1-p)$.
Gini: $G(p)=2p(1-p)$.

\textbf{Multiclass entropy} ($K$ classes, $p_k$=class fraction):
$H(Y)=-\sum_{k=1}^K p_k\log p_k$.

$IG(Y,X)=H(Y)-H(Y|X)$,\quad $H(Y|X)=\sum_v P(X{=}v)H(Y|X{=}v)$.

Real features: split at midpoints between adjacent values; $O(dm\log m)$ per level.

\textbf{Regularization}: max-depth, min-samples-split, min-impurity-decrease; or post-pruning (reduce-error pruning on val set).

VC: unrestricted trees $=\infty$; depth-$k$ tree VC $=O(2^k\log d)$.

Regression trees: minimize MSE within leaves; leaf value $=$ mean.


% =========================
\section*{Random Forest (Bagging)}

1) Bootstrap sample per tree.
2) Random subset $p'\!\ll\!d$ features at each split.
3) Predict: majority vote (classif.) / average (regr.).

Variance ($\rho$=pairwise tree correlation, $\sigma_f^2$=individual var.):
$\text{Var}(\bar f_M)=\rho\sigma_f^2+\frac{1-\rho}{M}\sigma_f^2$.

$M\!\to\!\infty$: variance $\to\rho\sigma_f^2$.

Small $p'$ $\Rightarrow$ lower $\rho$ $\Rightarrow$ lower ensemble variance.

\textbf{Key}: Bagging reduces variance NOT bias.


% =========================
\section*{AdaBoost}

Weak learner: $h_t\in\mathcal H$ with $\varepsilon_t<\tfrac12$ (better than random).

Init: $D_1(i)=1/m$. At round $t$:
1. Weighted err: $\varepsilon_t=\sum_{i:h_t(x_i)\ne y_i}D_t(i)$.
2. Weight: $\alpha_t=\tfrac12\log\frac{1-\varepsilon_t}{\varepsilon_t}$ ($>0$ since $\varepsilon_t<\tfrac12$).
3. Update: $D_{t+1}(i)=\frac{D_t(i)\exp(-\alpha_t y_i h_t(x_i))}{Z_t}$, $Z_t$ = normalizer.
   Misclassified points get higher weight; correctly classified lower.

Final: $H(x)=\text{sign}\!\left(\sum_t\alpha_t h_t(x)\right)$.

\textbf{Training error bound} ($\gamma_t=\tfrac12-\varepsilon_t>0$ = edge):
$\frac1m\sum_i\mathbf{1}[H(x_i)\ne y_i]\le\prod_t Z_t=e^{-2\sum_t\gamma_t^2}$.

Exponential loss connection: AdaBoost minimizes $\frac1m\sum_i e^{-y_i F(x_i)}$, $F=\sum_t\alpha_t h_t$.

Training error $\to0$ exponentially; generalization depends on margin.


% =========================
% GROUP 8: UNSUPERVISED
% =========================
\section*{PCA}

\textbf{Setup}: $X_c\in\mathbb R^{n\times d}$ (centered: subtract col-means).

Covariance: $S=\tfrac1n X_c^TX_c\in\mathbb R^{d\times d}$ (or $\tfrac1n XX^T$ with $X$ feature-row form).

Eigen-decomp: $S=V\Lambda V^T$, $\lambda_1\ge\lambda_2\ge\cdots\ge\lambda_d\ge0$.

PC scores: $Z=X_c V$ (project onto eigenvectors).

\textbf{EVR}: $\text{EVR}_j=\lambda_j/\sum_k\lambda_k$; $\text{CumEVR}_r=\sum_{j=1}^r\lambda_j/\sum_k\lambda_k$.

Pick $r$: CumEVR$\ge0.95$ or elbow in scree plot.

\textbf{Eckart-Young}: best rank-$r$ approx minimizes Frobenius error; reconstruction error $=\sum_{j>r}\lambda_j$.

\textbf{SVD link}: $X_c=U\Sigma V^T$ $\Rightarrow$ $\lambda_j=\sigma_j^2/n$.

\textbf{Whitening}: $\tilde z_j=z_j/\sqrt{\lambda_j}$ (unit variance per PC).

\textbf{Reconstruction}: $\hat x^i=Z_r V_r^T+\mathbf{1}\mu^T$ (add back means).

\textbf{Traps}: (1) must center first; (2) $n$ vs $n-1$; (3) EVR$\ne$eigenvalue; (4) SVD$\Sigma^2/n$ = eigenvalues; (5) $V$ cols not rows = PCs.


% =========================
\section*{Clustering (K-Means)}

\textbf{Objective}: $\min_{C_1,\ldots,C_K}\sum_{k=1}^K\sum_{x\in C_k}\|x-\mu_k\|^2$ (WCSS/SSD).

\textbf{Algorithm}: init $\mu_k$ $\to$ assign $x^i$ to nearest $\mu_k$ $\to$ recompute $\mu_k=\text{mean}(C_k)$ $\to$ repeat until stable.

WCSS non-increasing per iteration; terminates (finite partitions); finds local minimum.

\textbf{K-means++}: choose 1st centroid uniformly; each next $\propto d_i^2$ (min sq.\ dist.\ to nearest chosen). Gives $O(\log K)\times$OPT bound.

\textbf{TSS $=$ WCSS $+$ BCSS}: total $=$ within $+$ between cluster variance.

\textbf{Choosing $K$}:

\textit{Elbow}: plot WCSS vs $K$; pick knee point (diminishing returns).

\textit{Silhouette}: $a_i=$ mean intra-cluster dist; $b_i=$ min mean inter-cluster dist.
$s_i=\frac{b_i-a_i}{\max(a_i,b_i)}\in[-1,1]$; $s_D=\frac1n\sum_i s_i$. Higher = better.

\textit{Calinski-Harabasz}: $\text{CH}(K)=\frac{\text{BCSS}/(K-1)}{\text{WCSS}/(n-K)}$. Higher = better.

\textit{Davies-Bouldin}: $\text{DB}=\frac1K\sum_j\max_{k\ne j}\frac{s_j+s_k}{d_{jk}}$ ($s_j$=intra-scatter, $d_{jk}$=centroid dist). Lower = better.


% =========================
% GROUP 9: NEURAL NETWORKS
% =========================
\section*{MLP / Neural Networks}

\textbf{Architecture}: $L$ layers; $L_0$=input ($d$ units), $L_1,\ldots,L_{L-1}$=hidden, $L_L$=output.

Neuron: $z_j^\ell=\sum_k w_{jk}^\ell a_k^{\ell-1}$, $a_j^\ell=\phi(z_j^\ell)$.

Activations: sigmoid $\phi(z)=\tfrac{1}{1+e^{-z}}$, $\phi'=\phi(1-\phi)$; ReLU $\phi(z)=\max(0,z)$.

Motivation: XOR not linearly sep.; hidden layer maps to sep.\ space.

\textbf{Forward pass} (binary, 1 output):
$\hat y=\phi(W^L\phi(W^{L-1}\cdots\phi(W^1 x)\cdots))$.

Loss (MSE): $E=\sum_s({\hat y}^s-y^s)^2$.

\textbf{Backpropagation} (chain rule):
$\nabla_{W^\ell}e=\text{Diag}(\phi'^\ell)\,\delta^\ell\,(a^{\ell-1})^T$,

where $\delta^L=\frac{\partial e}{\partial a^L}$ and $\delta^\ell=V^{\ell+1}\cdots V^L\delta^L$,\quad $V^{\ell+1}=(W^{\ell+1})^T\text{Diag}(\phi'^{\ell+1})$.

Example: $\nabla_{w_{11}^3}e=\frac{\partial e}{\partial y}\phi'(z_1^3)a_1^2$.

\textbf{Optimization}: GD, SGD ($w^{k+1}\leftarrow w^k-\gamma_k\nabla_w e^{i_k}$), mini-batch SGD. Asymptotic convergence under $\sum\gamma_k^2<\infty$, $\sum\gamma_k=\infty$.


% =========================
\section*{MLP Multiclass}

Output layer: $C$ neurons; apply \textbf{softmax}:
$q_j=\frac{\exp(a_j^L)}{\sum_{r=1}^C\exp(a_r^L)}$,\quad $j=1,\ldots,C$.

$q=(q_1,\ldots,q_C)$ is a prob.\ distribution.

\textbf{One-hot encoding}: $y=c$ $\mapsto$ $y^{\text{oh}}\in\{0,1\}^C$ with $1$ at position $c$.

\textbf{KL divergence}: $KL(p\|q)=\sum_j p_j\log\frac{p_j}{q_j}=NE(p)+CE(p,q)$.

Since $NE(p)$ is constant, \textbf{minimizing KL $\equiv$ minimizing CE}.

\textbf{Cross-entropy loss}: $CE(p,q)=-\sum_j p_j\log q_j$.

With one-hot: $CE=-\log q_c$ (only true class contributes).

\textbf{Training}: forward pass $\to$ CE loss $\to$ backprop $\to$ SGD update.

Predict: $\hat y=\arg\max_j q_j$.

Backprop gradient: $\nabla_{W^\ell}e=\text{Diag}(\phi'^\ell)\delta^\ell(a^{\ell-1})^T$; $\delta^L=\partial e/\partial a^L$ differs for CE vs MSE.


% =========================
% GROUP 10: EXAM PREP
% =========================
\section*{Likely New Proof Questions}

\textbf{P1. K-means WCSS non-increasing.}
\textit{Assignment}: reassigning $x$ to nearest $\mu_k$ cannot increase $\|x-\mu_k\|^2$, so WCSS $\downarrow$ or stays.
\textit{Update}: $\min_\mu\sum_{x\in C_k}\|x-\mu\|^2$ $\Rightarrow$ $\partial/\partial\mu=0$ $\Rightarrow$ $\mu_k=\text{mean}(C_k)$; any other choice gives larger WCSS. Hence WCSS non-increasing per iteration.

\textbf{P2. TSS = WCSS + BCSS.}
$\text{TSS}=\sum_i\|x_i-\mu\|^2$. Write $x_i-\mu=(x_i-\mu_k)+(\mu_k-\mu)$; expand squared norm; cross term $\sum_{x\in C_k}\langle x_i-\mu_k,\mu_k-\mu\rangle=\langle\sum x_i-|C_k|\mu_k,\mu_k-\mu\rangle=0$ (since $\mu_k=\text{mean}$). Sum: $\text{TSS}=\text{WCSS}+\sum_k|C_k|\|\mu_k-\mu\|^2=\text{WCSS}+\text{BCSS}$.

\textbf{P3. PCA: PC1 = top eigenvector (Rayleigh quotient).}
$\max_{\|v\|=1}v^TSv$. Lagrangian: $Sv=\lambda v$, so $v$ is eigenvector. Objective $=v^TSv=\lambda$. Maximized by eigenvector with largest $\lambda$. \textit{Orthogonality}: subsequent PCs maximize variance subject to orthogonality to previous PCs $\Rightarrow$ $j$-th eigenvector.

\textbf{P4. Ridge = MAP with Gaussian prior.}
Prior: $p(w)\propto\exp(-\frac{\alpha}{2}\|w\|^2)$. Likelihood (Gaussian): $p(y|X,w)\propto\exp(-\frac{1}{2\sigma^2}\|y-Xw\|^2)$. MAP: $\arg\max_w[\log p(y|X,w)+\log p(w)]=\arg\min_w[\|y-Xw\|^2+\alpha\sigma^2\|w\|^2]$. Set $\lambda=\alpha\sigma^2$ $\Rightarrow$ Ridge.

\textbf{P5. LASSO = MAP with Laplace prior.}
Prior: $p(w_j)\propto\exp(-\lambda|w_j|)$, so $\log p(w)=-\lambda\|w\|_1$+const. MAP: $\arg\min_w[\|y-Xw\|^2+2\sigma^2\lambda\|w\|_1]$ $\Rightarrow$ LASSO. L1 penalty corners of L1 ball hit constraint surface at axes $\Rightarrow$ sparse solution.

\textbf{P6. Softmax sums to 1 and $q_j>0$.}
$q_j=\exp(a_j^L)/Z$, $Z=\sum_r\exp(a_r^L)>0$. Clearly $q_j>0$ (exp always positive). $\sum_jq_j=\sum_j\exp(a_j)/Z=Z/Z=1$. Hence valid probability distribution.

\textbf{P7. KL divergence non-negative.}
$KL(p\|q)=\sum_jp_j\log(p_j/q_j)\ge0$ by Jensen's inequality: $-KL=-\sum p_j\log(q_j/p_j)\ge-\log(\sum p_j\cdot q_j/p_j)=-\log 1=0$ ($\log$ concave). Equality iff $p=q$.

\textbf{P8. Bagging variance formula.}
Let $f_1,\ldots,f_M$ be trees with $\text{Var}(f_i)=\sigma_f^2$ and pairwise correlation $\rho$. $\text{Var}(\bar f_M)=\text{Var}(\frac1M\sum f_i)=\frac1{M^2}[M\sigma_f^2+M(M-1)\rho\sigma_f^2]=\frac{\sigma_f^2(1+(M-1)\rho)}{M}=\rho\sigma_f^2+\frac{(1-\rho)\sigma_f^2}{M}$. As $M\to\infty$: $\to\rho\sigma_f^2$ (irreducible if trees correlated).

\textbf{P9. CE loss is convex in logits (softmax output layer).}
For binary: $\ell(w)=-y\log\sigma(z)-(1-y)\log(1-\sigma(z))$. Second deriv $=\sigma(z)(1-\sigma(z))>0$ $\Rightarrow$ convex in $z=w^Tx$ $\Rightarrow$ convex in $w$ (affine composition). For multiclass CE with softmax: Hessian $H=\text{diag}(q)-qq^T\succeq0$ (Schur complement argument).

\textbf{P10. Why L1 gives sparsity, L2 does not.}
Ridge: constraint $\|w\|^2\le t$ (sphere); loss contours ellipsoidal; intersection generically at non-axis point $\Rightarrow$ no sparsity. LASSO: constraint $\|w\|_1\le t$ (diamond); corners lie on coordinate axes; loss contours likely touch corner first $\Rightarrow$ many $w_j=0$. In 2D: corners at $(t,0),(0,t),(-t,0),(0,-t)$.


% =========================
\section*{More Computation Questions}

\textbf{Q: Perceptron convergence bound numeric.}
Given $\|w^*\|=2$, $\rho=3$, $\gamma=0.5$: $T\le(\rho/\gamma)^2=(3/0.5)^2=36$. Also $T\le\rho^2\|w^*\|^2/\gamma^2=4\times9/0.25=144$ (multiclass version). Report bound from margin formula.

\textbf{Q: Naive Bayes prediction.}
Given class priors $P(y)$ and likelihoods $P(x_j|y)$: compute $\log P(y)+\sum_j\log P(x_j|y)$ for each class; predict class with max log-score. Apply Laplace smoothing if any count is zero: add $\alpha=1$ to numerator, $\alpha|\mathcal V|$ to denominator.

\textbf{Q: Compute F1 from raw counts.}
Given TP=50, FP=10, FN=20: Precision$=50/60\approx0.83$; Recall$=50/70\approx0.71$; $F1=2\times0.83\times0.71/(0.83+0.71)\approx0.77$. Note: F1 harmonic mean of P and R, not arithmetic.

\textbf{Q: Random Forest variance reduction.}
Single tree: $\sigma_f^2=4$. $M=100$ trees, pairwise correlation $\rho=0.1$. $\text{Var}(\bar f)=0.1\times4+\frac{0.9\times4}{100}=0.4+0.036=0.436$. Single tree variance = 4; RF reduces to 0.436 ($\approx10\times$ reduction). Floor is $\rho\sigma_f^2=0.4$.

\textbf{Q: Softmax + CE by hand.}
Logits $a=[2,1,0]$. $Z=e^2+e^1+e^0=7.389+2.718+1=11.107$. $q=[0.665,0.245,0.090]$. True class $c=1$: $CE=-\log(0.665)=0.408$. Gradient: $\nabla_{a}CE=q-e_c=[0.665-1,0.245,0.090]=[-0.335,0.245,0.090]$.


% =========================
\section*{More Proof Questions}

\textbf{P11. AdaBoost training error exponentially decays.}
Show $\frac1m\sum_i\mathbf1[H(x_i)\ne y_i]\le e^{-2\sum_t\gamma_t^2}$.
Key: $\mathbf1[H(x_i)\ne y_i]\le e^{-y_iF(x_i)}$ (since $y_iF<0\Rightarrow e^{-y_iF}\ge1$). Sum: $\frac1m\sum_i e^{-y_iF(x_i)}$. Show inductively this equals $\prod_tZ_t$: at each round, $D_{t+1}$ is renormalized, and $Z_t=\sqrt{4\varepsilon_t(1-\varepsilon_t)}=\sqrt{1-4\gamma_t^2}\le e^{-2\gamma_t^2}$.

\textbf{P12. K-means always terminates.}
Number of distinct partitions of $n$ points into $K$ clusters is finite ($K^n$). Each iteration of K-means either decreases WCSS strictly or keeps it same (only if assignment unchanged). Since WCSS is bounded below by 0 and decreases between distinct partitions, algorithm must terminate at a fixed point.

\textbf{P13. Entropy $H(p)$ maximized by uniform distribution.}
Maximize $-\sum_{k=1}^K p_k\log p_k$ subject to $\sum p_k=1$, $p_k\ge0$. Lagrangian: $\partial/\partial p_k[{-p_k\log p_k-\lambda p_k}]=0\Rightarrow-\log p_k-1-\lambda=0\Rightarrow p_k=e^{-1-\lambda}=$ constant. Hence $p_k=1/K$ for all $k$; max entropy $=\log K$.

\textbf{P14. Ridge always has unique solution.}
$\hat\beta=(X^TX+\lambda I)^{-1}X^Ty$ for any $\lambda>0$. Show $X^TX+\lambda I\succ0$: for any $v\ne0$, $v^T(X^TX+\lambda I)v=\|Xv\|^2+\lambda\|v\|^2\ge\lambda\|v\|^2>0$. Hence positive definite $\Rightarrow$ invertible $\Rightarrow$ unique solution, even when $X^TX$ is singular (collinear features).

\textbf{P15. VC dimension of linear classifiers in $\mathbb R^d$ is $d+1$.}
\textit{Achievability}: $d+1$ points in general position (no $d+1$ on a hyperplane) are shattered by homogeneous halfspaces in $\mathbb R^{d+1}$. \textit{Upper bound}: any $d+2$ points are linearly dependent; any labeling cannot be realized. Formally, Radon's theorem: any $d+2$ points can be split into two sets with convex hulls intersecting $\Rightarrow$ no classifier separates all $2^{d+2}$ labelings.

\section*{HIGH PROBABILITY QUESTIONS --- SOLUTIONS}

\textbf{PAPER SET A SOLUTIONS.}

\textbf{Short Answers.}
\textbf{1. Bias-Variance.} Increase $\lambda\Rightarrow$ bias$\uparrow$, variance$\downarrow$, overfitting$\downarrow$, underfitting risk$\uparrow$.

\textbf{2. Bayes classifier.} $h^*(x)=\mathbf{1}[\eta(x)>1/2]$, $\eta(x)=P(Y=1|x)$.

\textbf{3. Support vectors.} Points with $y_i(w^Tx_i+b)=1$ (hard margin); they determine the boundary.

\textbf{4. Generative vs Discriminative.} Generative: model $P(X,Y)$ / $P(X|Y),P(Y)$ (NB). Discriminative: model $P(Y|X)$ or boundary (Logistic, SVM).

\textbf{5. Scaling in KNN.} Large-scale features dominate Euclidean distance.

\textbf{6. Hinge loss.} $\max(0,1-yf(x))$; convex surrogate of $0$--$1$.

\textbf{7. $h_{ii}$.} Leverage of point $i$; large $h_{ii}\Rightarrow$ influential.

\textbf{8. Bagging vs Boosting.} Bagging reduces variance. Boosting reduces bias / builds strong learner.

\textbf{9. PCA centering.} Need covariance about mean; otherwise PC1 is pulled toward mean vector.

\textbf{10. Slack variables.} $\xi_i$ allow margin violations / misclassification.

\textbf{Medium Answers.}
\textbf{11. Logistic gradient.} $J=-\sum[y\log p+(1-y)\log(1-p)]$, $p=\sigma(Xw)$, $\nabla J=X^T(p-y)$.

\textbf{12. Perceptron updates.} $w_0=(0,0)$. $x_1=(2,1),y=+1$ mistake $\Rightarrow w_1=(2,1)$. $x_2=(-1,-1),y=-1$: $y(w^Tx)=(-1)(-3)=3>0$ correct. Final $w=(2,1)$.

\textbf{13. Ridge.} $\min \|y-X\beta\|^2+\lambda\|\beta\|^2$. Set grad $=0$:
\[
(X^TX+\lambda I)\beta=X^Ty,\qquad \hat\beta=(X^TX+\lambda I)^{-1}X^Ty.
\]
Invertible since $v^T(X^TX+\lambda I)v=\|Xv\|^2+\lambda\|v\|^2>0$ for $v\neq0$.

\textbf{14. PCA.} $|S-\lambda I|=(4-\lambda)(3-\lambda)-4=\lambda^2-7\lambda+8=0$, so
\[
\lambda_{1,2}=\frac{7\pm\sqrt{17}}{2},\qquad \lambda_{\max}=\frac{7+\sqrt{17}}{2}.
\]
Eigenvector: $(4-\lambda)x+2y=0\Rightarrow y=\frac{\lambda-4}{2}x$, use any proportional vector.

\textbf{15. K-means.} $C_1=\{(1,1),(2,1)\}\Rightarrow \mu_1=(1.5,1)$. $C_2=\{(8,8),(9,8)\}\Rightarrow \mu_2=(8.5,8)$.

\textbf{Long Proof Skeletons.}
\textbf{16. Bayes optimality.} At fixed $x$: $\mathrm{Err}(g|x)=1-P(g(x)|x)$. Minimize by choosing $\arg\max_y P(y|x)=h^*(x)$. Then average over $x$: $R(h^*)\le R(g)$.

\textbf{17. Perceptron.} $\langle w^*,w_T\rangle\ge T\gamma$, $\|w_T\|^2\le T\rho^2$. By Cauchy:
\[
T\gamma\le \|w^*\|\|w_T\|\le \|w^*\|\rho\sqrt T
\Rightarrow T\le \frac{\rho^2\|w^*\|^2}{\gamma^2}.
\]
If $\|w^*\|=1$, $T\le(\rho/\gamma)^2$.

\textbf{18. Hard SVM dual.} Primal $\min \frac12\|w\|^2$ s.t. $y_i(w^Tx_i+b)\ge1$. Lagrangian
\[
L=\frac12\|w\|^2-\sum_i\alpha_i[y_i(w^Tx_i+b)-1].
\]
Stationarity:
\[
\partial_wL=0\Rightarrow w=\sum_i\alpha_i y_i x_i,\qquad
\partial_bL=0\Rightarrow \sum_i\alpha_i y_i=0.
\]
Dual:
\[
\max_{\alpha\ge0}\sum_i\alpha_i-\frac12\sum_{i,j}\alpha_i\alpha_j y_iy_j x_i^Tx_j,
\quad \sum_i\alpha_i y_i=0.
\]

\textbf{PAPER SET B SOLUTIONS.}

\textbf{Short Answers.}
\textbf{1. OLS vs Ridge vs LASSO.} OLS: no penalty. Ridge: $\lambda\|\beta\|^2$ shrinkage. LASSO: $\lambda\|\beta\|_1$ sparsity.

\textbf{2. Precision/Recall/F1.}
\[
P=\frac{TP}{TP+FP},\quad R=\frac{TP}{TP+FN},\quad F1=\frac{2PR}{P+R}.
\]

\textbf{3. Logistic no closed form.} Normal equations become nonlinear due sigmoid.

\textbf{4. Overfitting.} Low train error, high test error.

\textbf{5. Trees overfit.} Deep splits can memorize noise.

\textbf{6. Curse of dimensionality.} Distances concentrate; NN becomes unreliable.

\textbf{7. Kernel trick.} Replace $x^Tz$ by $K(x,z)=\phi(x)^T\phi(z)$.

\textbf{8. LOOCV.} Train $n$ times, leave one point out each time.

\textbf{9. Bagging variance.} Averaging decorrelated learners lowers variance.

\textbf{10. Softmax.}
\[
q_j=\frac{e^{a_j}}{\sum_k e^{a_k}}.
\]

\textbf{Medium.}
\textbf{11. TP50 FP10 FN20.}
\[
P=\frac{50}{60}=0.833,\qquad R=\frac{50}{70}=0.714,\qquad
F1\approx0.769.
\]

\textbf{12. Softmax$(2,1,0)$.}
\[
Z=e^2+e+1=11.107,\quad q=(0.665,0.245,0.090),
\]
CE for true class 1st/index 1:
\[
CE=-\log(0.665)=0.408.
\]

\textbf{13. LOOCV.}
\[
\hat y_i^{(-i)}=\hat y_i-\frac{h_{ii}e_i}{1-h_{ii}},\qquad
e_i^{(-i)}=\frac{e_i}{1-h_{ii}},
\]
\[
CV=\frac1n\sum_i\left(\frac{e_i}{1-h_{ii}}\right)^2.
\]

\textbf{14. Naive Bayes.}
\[
\text{score}(y)=\log P(y)+\sum_j\log P(x_j|y),\qquad \hat y=\arg\max_y \text{score}(y).
\]

\textbf{15. AdaBoost.}
\[
\epsilon_t=\sum_{i:\text{wrong}}D_t(i),\qquad
\alpha_t=\frac12\log\frac{1-\epsilon_t}{\epsilon_t}.
\]

\textbf{Long.}
\textbf{16. Hinge convex.} $\max\{0,1-yw^Tx\}$ is max of affine functions $\Rightarrow$ convex.

\textbf{17. Logistic convex.}
\[
\ell(z)=\log(1+e^{-z}),\qquad
\ell''(z)=\frac{e^z}{(1+e^z)^2}\ge0.
\]

\textbf{18. Ridge = MAP.} Gaussian prior $p(w)\propto e^{-\alpha\|w\|^2/2}$ + Gaussian likelihood $\Rightarrow$ MAP minimizes SSE$+\lambda\|w\|^2$.

\textbf{19. LASSO = MAP.} Laplace prior $p(w_j)\propto e^{-\lambda|w_j|}$ $\Rightarrow$ MAP minimizes SSE$+\lambda\|w\|_1$.

\textbf{20. 1NN.}
\[
R_{1NN}(x)=2\eta(x)(1-\eta(x))\le 2\min(\eta,1-\eta)=2R^*(x).
\]
Take expectation: $R_{1NN}\le2R^*$.

\textbf{PAPER SET C SOLUTIONS.}

\textbf{1. Representer.} Write $w=u+v$, $u\in\mathrm{span}\{x_i\}$, $v\perp \mathrm{span}\{x_i\}$. Loss depends only on $u$ since $v^Tx_i=0$. Norm: $\|w\|^2=\|u\|^2+\|v\|^2\ge\|u\|^2$. So optimum has $v=0$.

\textbf{2. Strong convexity.} If $w_1\neq w_2$ are minimizers and $f$ strongly convex, midpoint satisfies
\[
f\Big(\frac{w_1+w_2}{2}\Big)<\frac{f(w_1)+f(w_2)}{2},
\]
contradicting minimality.

\textbf{3. K-means monotonic.} Assignment step lowers WCSS by nearest-centroid choice. Update step lowers WCSS since mean minimizes within-cluster SSE. Therefore WCSS non-increasing.

\textbf{4. TSS decomposition.}
\[
\sum\|x-\mu\|^2=\sum\|x-\mu_k\|^2+\sum n_k\|\mu_k-\mu\|^2
\]
since cross term vanishes by $\sum_{x\in C_k}(x-\mu_k)=0$.

\textbf{5. PCA.}
\[
\max_{\|v\|=1}v^TSv
\]
gives Lagrangian condition $Sv=\lambda v$; largest $\lambda$ gives PC1.

\textbf{6. Weak duality.} For all $x,y$,
\[
\min_x f(x,y)\le f(x,y)\le \max_y f(x,y).
\]
Then maximize left side over $y$, minimize right side over $x$:
\[
\max_y\min_x f(x,y)\le \min_x\max_y f(x,y).
\]

\textbf{7. Hinge upper-bounds $0$--$1$.} If $yf\le0$, then $1-yf\ge1$, so hinge $\ge1=\mathbf{1}[yf\le0]$. If $yf>0$, both nonnegative with hinge $\ge0$.

\textbf{8. Logistic upper bound.} If $yf\le0$, then $-yf\ge0$ and
\[
\log(1+e^{-yf})\ge \log(1+1)=\log2.
\]

\textbf{9. Softmax valid.} $q_j>0$ since exponentials are positive and
\[
\sum_j q_j=\frac{\sum_j e^{a_j}}{\sum_k e^{a_k}}=1.
\]

\textbf{10. Bagging variance.}
\[
\mathrm{Var}(\bar f_M)=\rho\sigma^2+\frac{1-\rho}{M}\sigma^2.
\]

\textbf{LAST NIGHT NUMERICALS.}
\[
\hat\beta_{\text{ridge}}=(X^TX+\lambda I)^{-1}X^Ty,\qquad
\nabla J_{\text{logistic}}=X^T(p-y),
\]
\[
\gamma_{\text{SVM}}=\frac{1}{\|w\|},\qquad
T_{\text{perc}}\le(\rho/\gamma)^2,
\]
\[
LOOCV=\frac1n\sum_i\left(\frac{e_i}{1-h_{ii}}\right)^2,\qquad
F1=\frac{2TP}{2TP+FP+FN},
\]
\[
CE_{\text{softmax}}=-\log q_{\text{true}},\qquad
\mathrm{EVR}_j=\frac{\lambda_j}{\sum_r\lambda_r}.
\]

\section*{MOST PROBABLE ENDSEM PROOF QUESTIONS (HIGH YIELD)}

\textbf{TIER A : VERY HIGH PROBABILITY (Almost Certain).}

\textbf{1. Bayes Optimality.}
For any classifier $g$:
\[
R(h^*)\le R(g),\qquad h^*(x)=\arg\max_y P(Y=y|X=x).
\]
Proof: at fixed $x$,
\[
\mathrm{Err}(g|x)=1-P(g(x)|x).
\]
Bayes picks max posterior $\Rightarrow$ minimum conditional error. Then take expectation over $x$.

\textbf{2. Perceptron Mistake Bound.}
If separable with margin $\gamma$ and $\|x_i\|\le\rho$:
\[
T\le(\rho/\gamma)^2.
\]
Skeleton:
\[
\langle w^*,w_T\rangle\ge T\gamma,\qquad \|w_T\|^2\le T\rho^2.
\]
By Cauchy:
\[
T\gamma\le \|w^*\|\rho\sqrt T \Rightarrow T\le(\rho/\gamma)^2
\]
for normalized $\|w^*\|=1$.

\textbf{3. Hard Margin SVM = Max Margin.}
\[
\max_{\|w\|=1,b}\min_i y_i(w^Tx_i+b)
\]
is equivalent to
\[
\min_{w,b}\frac12\|w\|^2\quad\text{s.t. }y_i(w^Tx_i+b)\ge1.
\]
Reason: separator is scale-invariant. Fix min functional margin $=1$. Then geometric margin $=1/\|w\|$, so maximizing margin $\Leftrightarrow$ minimizing $\|w\|$.

\textbf{4. Dual of Hard SVM.}
Primal:
\[
\min \frac12\|w\|^2 \quad \text{s.t. } y_i(w^Tx_i+b)\ge1.
\]
Lagrangian:
\[
L=\frac12\|w\|^2-\sum_i\alpha_i[y_i(w^Tx_i+b)-1].
\]
Stationarity:
\[
w=\sum_i\alpha_i y_i x_i,\qquad \sum_i\alpha_i y_i=0.
\]
Dual:
\[
\max_\alpha \sum_i\alpha_i-\frac12\sum_{i,j}\alpha_i\alpha_j y_i y_j x_i^Tx_j
\]
s.t. $\alpha_i\ge0$, $\sum_i\alpha_i y_i=0$.

\textbf{5. Logistic Loss Convexity.}
\[
\ell(z)=\log(1+e^{-z}),\qquad \ell''(z)=\frac{e^z}{(1+e^z)^2}\ge0.
\]
Hence convex.

\textbf{TIER B : HIGH PROBABILITY.}

\textbf{6. Hinge Loss Convexity.}
\[
\ell(w)=\max(0,1-yw^Tx).
\]
Max of affine functions $\Rightarrow$ convex.

\textbf{7. Ridge Always Unique.}
For $\lambda>0$,
\[
v^T(X^TX+\lambda I)v=\|Xv\|^2+\lambda\|v\|^2>0 \quad (v\neq0).
\]
So $X^TX+\lambda I\succ0$, invertible, unique minimizer.

\textbf{8. Bias Variance Decomposition.}
\[
\mathbb{E}[(y-\hat f)^2]=(\mathbb{E}[\hat f]-f)^2+\mathrm{Var}(\hat f)+\sigma^2.
\]
Use $y=f+\epsilon$, $\mathbb{E}[\epsilon]=0$, expand square, kill cross terms.

\textbf{9. PCA First Principal Component.}
\[
\max_{\|v\|=1}v^TSv.
\]
Lagrangian:
\[
\mathcal L(v,\lambda)=v^TSv-\lambda(v^Tv-1).
\]
Stationarity gives
\[
Sv=\lambda v.
\]
Largest eigenvalue $\lambda$ gives PC1.

\textbf{10. K-Means Objective Never Increases.}
Assignment step: nearest centroid lowers WCSS. Update step: mean minimizes $\sum\|x-\mu\|^2$. Hence objective non-increasing each iteration.

\textbf{TIER C : MODERATE BUT POSSIBLE.}

\textbf{11. Representer Theorem (Linear).}
For
\[
\min \lambda\|w\|^2+\sum_i \ell(y_i,w^Tx_i),
\]
show $w^*=\sum_i\alpha_i x_i$.
Write $w=u+v$ with $u\in\mathrm{span}\{x_i\}$, $v\perp\mathrm{span}\{x_i\}$. Loss depends only on $u$ since $v^Tx_i=0$. Norm increases if $v\neq0$. Hence $v=0$.

\textbf{12. Strong Convexity $\Rightarrow$ Unique Minimizer.}
If $f$ is $\lambda$-strongly convex, two distinct minimizers $w_1,w_2$ impossible because midpoint gives strictly lower value.

\textbf{13. Weak Duality.}
\[
\min_x\max_y f(x,y)\ge \max_y\min_x f(x,y).
\]
Since for all $x,y$:
\[
\min_x f(x,y)\le f(x,y)\le \max_y f(x,y).
\]

\textbf{14. Hinge Upper Bounds 0--1.}
\[
\mathbf{1}[yf\le0]\le \max(0,1-yf).
\]

\textbf{15. Logistic Upper Bounds 0--1.}
\[
\mathbf{1}[yf\le0]\le \log(1+e^{-yf})
\]
up to constant / $\log2$ scaling idea.

\textbf{TIER D : If KNN/Bayes Asked.}

\textbf{16. 1-NN Risk Bound.}
\[
R_{1NN}\le 2R^*.
\]
At fixed $x$, $\eta=P(Y=1|x)$,
\[
R_{1NN}(x)=2\eta(1-\eta)\le2\min(\eta,1-\eta)=2R^*(x).
\]
Take expectation.

\textbf{ULTRA SHORT MEMORIZE ORDER (If no time).}
1. Bayes optimality
2. Perceptron bound
3. Hard SVM primal-dual
4. Logistic convexity
5. Hinge convexity
6. PCA eigenvector proof
7. Kmeans monotonicity
8. Ridge invertibility
9. Bias-variance
10. $1$NN $\le 2R^*$

\textbf{EXAM STRATEGY.}
If proof unknown:
1. Write definitions
2. Use convexity / Cauchy / Lagrangian / expectation
3. Show final theorem statement
4. Even partial derivation gets marks

\textbf{VERY LIKELY NUMERICAL + PROOF MIXED.}
\begin{itemize}
\item Find support vectors then explain KKT.
\item PCA eigenvalues then justify PC1.
\item Kmeans one iteration then prove WCSS decreases.
\item Ridge solution then prove invertibility.
\item Logistic gradient then prove convexity.
\end{itemize}

\textbf{EXTRA PROOF-TYPE QUESTIONS + SOLNS.}

\textbf{17. OLS normal equations.} Show minimizer of $\|y-X\beta\|^2$ satisfies $X^TX\beta=X^Ty$. Solution: $\nabla_\beta\|y-X\beta\|^2=-2X^T(y-X\beta)=0$.

\textbf{18. Hat matrix symmetric/idempotent.} $H=X(X^TX)^{-1}X^T$. Then $H^T=H$ and $H^2=X(X^TX)^{-1}(X^TX)(X^TX)^{-1}X^T=H$.

\textbf{19. Residuals orthogonal to column space.} $e=y-\hat y=(I-H)y$. Then $X^Te=X^Ty-X^TX\hat\beta=0$; also $He=H(I-H)y=0$.

\textbf{20. Information gain nonnegative.} $IG(Y;X)=H(Y)-H(Y|X)=I(X;Y)\ge0$ since conditioning cannot increase entropy / KL$\ge0$.

\textbf{21. Covariance matrix PSD.} For covariance $S$, $v^TSv=\mathrm{Var}(v^TX)\ge0$ for all $v$, hence $S\succeq0$.

\textbf{22. PCA eigenvectors orthogonal.} If $Su=\lambda u$, $Sv=\mu v$, $\lambda\neq\mu$, then $u^TSv=\mu u^Tv=(Su)^Tv=\lambda u^Tv$, so $(\lambda-\mu)u^Tv=0$.

\textbf{23. Softmax shift invariance.} $\displaystyle q_j(a+c\mathbf1)=\frac{e^{a_j+c}}{\sum_k e^{a_k+c}}=\frac{e^{a_j}}{\sum_k e^{a_k}}=q_j(a)$.

\textbf{24. Cross-entropy minimized at truth.} $CE(p,q)=H(p)+KL(p\|q)$. Since $KL\ge0$, minimum occurs at $q=p$.

\textbf{25. Bagging variance reduction.} $\mathrm{Var}(\bar f_M)=\rho\sigma^2+\frac{1-\rho}{M}\sigma^2\le\sigma^2$ for $\rho\le1$; strict drop if $\rho<1$ and $M>1$.

\textbf{26. Tree training error non-increasing with depth.} A deeper tree can copy parent prediction by making no effective split; best split cannot worsen empirical impurity/error, so train error never increases.

\textbf{27. LOOCV shortcut.} With hat matrix $H$, $e_i^{(-i)}=\frac{e_i}{1-h_{ii}}$. Follows from Sherman--Morrison update after deleting row $i$.

\textbf{28. Projection non-expansive.} For convex set $C$, $\|\Pi_C(u)-\Pi_C(v)\|\le\|u-v\|$. Use projection optimality: $\langle u-\Pi_C(u),z-\Pi_C(u)\rangle\le0$ with $z=\Pi_C(v)$ and swap $u,v$.

\textbf{29. Cluster mean minimizes SSE.} $\arg\min_\mu\sum_{i\in C}\|x_i-\mu\|^2$ satisfies $\partial/\partial\mu=-2\sum_{i\in C}(x_i-\mu)=0$, so $\mu=\frac1{|C|}\sum_{i\in C}x_i$.

\textbf{30. TSS = WCSS + BCSS.} Write $x_i-\mu=(x_i-\mu_k)+(\mu_k-\mu)$; square and sum. Cross term vanishes since $\sum_{i\in C_k}(x_i-\mu_k)=0$.

\textbf{31. PCs are uncorrelated.} If $Z=XV$ with $S=V\Lambda V^T$, then $\mathrm{Cov}(Z)=V^TSV=\Lambda$, diagonal. Hence PC coordinates are uncorrelated.

\textbf{32. OLS unbiased.} With $y=X\beta+\epsilon$, $\mathbb E[\epsilon|X]=0$, $\hat\beta=(X^TX)^{-1}X^Ty=\beta+(X^TX)^{-1}X^T\epsilon$, so $\mathbb E[\hat\beta|X]=\beta$.

\textbf{33. Ridge biased.} $\hat\beta_R=(X^TX+\lambda I)^{-1}X^Ty$, so $\mathbb E[\hat\beta_R|X]=(X^TX+\lambda I)^{-1}X^TX\beta\neq\beta$ for $\lambda>0$.

\textbf{34. Ridge shrinks singular directions.} If $X=UDV^T$, then $\hat\beta_R=V(D^2+\lambda I)^{-1}DU^Ty$; coefficient on singular direction $j$ gets factor $\frac{d_j}{d_j^2+\lambda}$.

\textbf{35. Bernoulli NB gives linear log-odds.} $\log\frac{P(y=1|x)}{P(y=0|x)}=\log\frac{\pi}{1-\pi}+\sum_j x_j\log\frac{\theta_{j1}(1-\theta_{j0})}{\theta_{j0}(1-\theta_{j1})}+\text{const}$, linear in $x$.

\textbf{36. Gaussian equal covariance $\Rightarrow$ linear boundary.} Compare class discriminants $-\frac12(x-\mu_k)^T\Sigma^{-1}(x-\mu_k)+\log\pi_k$; quadratic $x^T\Sigma^{-1}x$ cancels when $\Sigma$ same for both classes.

\textbf{37. Unequal covariance $\Rightarrow$ quadratic boundary.} If $\Sigma_1\neq\Sigma_0$, the $x^T\Sigma_k^{-1}x$ terms do not cancel, so decision boundary is quadratic.

\textbf{38. Logistic Hessian PSD.} For binary logistic, $\nabla^2J=X^TWX$ with $W=\mathrm{diag}(p_i(1-p_i))\succeq0$; hence $v^T\nabla^2Jv=(Xv)^TW(Xv)\ge0$.

\textbf{39. Valid kernel $\Rightarrow$ Gram PSD.} If $K(x,z)=\phi(x)^T\phi(z)$, then for any $c$, $c^TKc=\|\sum_i c_i\phi(x_i)\|^2\ge0$.

\textbf{40. AdaBoost upweights mistakes.} $D_{t+1}(i)\propto D_t(i)e^{-\alpha_ty_ih_t(x_i)}$: wrong point has $y_ih_t=-1\Rightarrow e^{+\alpha_t}$; correct point has factor $e^{-\alpha_t}$.

\textbf{41. AdaBoost $\alpha_t$ formula.} Minimize $Z_t(\alpha)=\epsilon_te^\alpha+(1-\epsilon_t)e^{-\alpha}$; set derivative $0$ to get $\alpha_t=\frac12\log\frac{1-\epsilon_t}{\epsilon_t}$.

\textbf{42. Geometric margin scale-invariant.} $\gamma_i=\frac{y_i(w^Tx_i+b)}{\|w\|}$ is unchanged by $(w,b)\mapsto(cw,cb)$ for $c>0$, since numerator and denominator both scale by $c$.

\textbf{43. Non-support points do not matter.} In dual, $w=\sum_i\alpha_i y_i x_i$; if $\alpha_i=0$ the point contributes nothing, so only support vectors determine $w$.

\textbf{44. Soft-SVM: $\alpha_i=C$ means inside margin/violation.} Since $\alpha_i+\beta_i=C$, $\alpha_i=C\Rightarrow\beta_i=0$; KKT allows $\xi_i>0$, so point is on/inside wrong side of margin.

\textbf{45. Hard-SVM support vectors lie on margin.} If $\alpha_i>0$, complementary slackness gives $\alpha_i[y_i(w^Tx_i+b)-1]=0$, so $y_i(w^Tx_i+b)=1$.

\textbf{46. Trace of hat matrix.} $\mathrm{tr}(H)=\mathrm{tr}(X(X^TX)^{-1}X^T)=\mathrm{tr}((X^TX)^{-1}X^TX)=p$ for full-rank $X\in\mathbb R^{n\times p}$.

\textbf{47. Residual sum zero with intercept.} If first column of $X$ is $\mathbf 1$, then $\mathbf1^Te=0$ since residuals are orthogonal to every column of $X$.

\textbf{48. $R^2$ identity.} With intercept, $R^2=\frac{SSM}{SST}=1-\frac{SSE}{SST}$ because $SST=SSM+SSE$.

\textbf{49. PCA needs centering.} Without centering, covariance gets contaminated by mean term $\mu\mu^T$; PC1 can point toward mean instead of variation.

\textbf{50. Softmax-CE gradient.} For one-hot target $y$, $\nabla_a CE=q-y$. Differentiate $-\sum_j y_j\log q_j$ with $q_j=e^{a_j}/\sum_k e^{a_k}$.

\textbf{51. Logistic decision boundary linear.} Predict class $1$ when $\sigma(w^Tx+b)>\frac12$, equivalent to $w^Tx+b>0$, a hyperplane.

\textbf{52. Exponential upper-bounds $0$--$1$.} $\mathbf1[yf\le0]\le e^{-yf}$ since if $yf\le0$ then $e^{-yf}\ge1$, else RHS $>0$.

\textbf{53. Exponential loss convex.} $e^{-z}$ has second derivative $e^{-z}>0$, so convex; composition with affine $z=yw^Tx$ stays convex.

\textbf{54. AdaBoost train-error bound.} $\frac1m\sum_i\mathbf1[H(x_i)\neq y_i]\le\frac1m\sum_i e^{-y_iF(x_i)}=\prod_t Z_t$, and $Z_t\le e^{-2\gamma_t^2}$.

\textbf{55. Entropy maximized at uniform.} Lagrangian for $-\sum_k p_k\log p_k$ under $\sum_k p_k=1$ gives all $p_k$ equal, so $p_k=1/K$.

\textbf{56. Mutual information symmetric.} $I(X;Y)=H(X)-H(X|Y)=H(Y)-H(Y|X)=H(X)+H(Y)-H(X,Y)$.

\textbf{57. One-hot CE = NLL.} If true class is $c$, $CE=-\sum_j y_j\log q_j=-\log q_c$, same as negative log-likelihood.

\textbf{58. Random-forest variance limit.} As $M\to\infty$, $\mathrm{Var}(\bar f_M)\to\rho\sigma^2$; averaging removes only uncorrelated part.

\textbf{59. LASSO KKT sparsity test.} At optimum, if $\beta_j=0$ then $|x_j^T(y-X\beta)|\le\lambda$; if $|\cdot|<\lambda$, coordinate stays exactly zero.

\textbf{60. Orthogonal-design LASSO = soft threshold.} If $X^TX=I$, then $\hat\beta_j=\mathrm{sign}(z_j)(|z_j|-\lambda)_+$ with $z_j=x_j^Ty$.

\textbf{61. K-means terminates finitely.} Finite number of assignments; each changed assignment lowers WCSS; bounded below by $0$, so algorithm stops at local optimum.

\textbf{62. Conditional mean minimizes MSE.} Minimize $\mathbb E[(Y-a)^2|X=x]$; derivative $-2\mathbb E[Y-a|x]=0$ gives $a=\mathbb E[Y|X=x]$.

\textbf{63. Conditional median minimizes MAE.} Minimize $\mathbb E[|Y-a||X=x]$; any median makes left/right subgradient probabilities balance.

\textbf{64. Conditional mode minimizes $0$--$1$.} For classifier predicting label $a$, conditional risk is $1-P(Y=a|x)$, minimized by $a=\arg\max_y P(Y=y|x)$.

\textbf{65. Cross-entropy lower bound.} $CE(p,q)=H(p)+KL(p\|q)\ge H(p)$.

\textbf{66. KL zero iff equal distributions.} $KL(p\|q)\ge0$ with equality iff $p_j=q_j$ for all $j$ with $p_j>0$.

\textbf{67. Bayes threshold with unequal costs.} Predict $1$ iff $C_{10}P(Y=0|x)\le C_{01}P(Y=1|x)$, i.e. $\eta(x)\ge \frac{C_{10}}{C_{01}+C_{10}}$.

\textbf{68. Gaussian NB means diagonal covariance.} Conditional independence of features under class $y$ implies $\Sigma_y$ has zero off-diagonals, so class density factorizes.

\textbf{69. Kernel SVM prediction.} Since $w=\sum_i\alpha_i y_i\phi(x_i)$, $f(x)=w^T\phi(x)+b=\sum_i\alpha_i y_i K(x_i,x)+b$.

\textbf{70. Nearest-centroid boundary linear.} Compare $\|x-\mu_1\|^2$ and $\|x-\mu_2\|^2$; $x^Tx$ cancels, leaving linear boundary.

\textbf{71. Perceptron raises mistaken score.} After $w^+=w+yx$, $y\langle w^+,x\rangle=y\langle w,x\rangle+\|x\|^2$.

\textbf{72. Hinge subgradient.} $0$ if $yw^Tx>1$, $-yx$ if $yw^Tx<1$, any convex combo at equality.

\textbf{73. Logistic derivative.} $\sigma'(z)=\sigma(z)(1-\sigma(z))$.

\textbf{74. PCA variance along $v_j$.} $\mathrm{Var}(v_j^Tx)=v_j^TSv_j=\lambda_j$.

\textbf{75. EVRs sum to 1.} $\sum_j \lambda_j/\sum_r\lambda_r=1$.

\textbf{76. Best tree split non-worsening.} Chosen split minimizes weighted child impurity, so post-split impurity $\le$ parent impurity.

\textbf{77. Leaf majority vote optimal.} In a leaf, class minimizing $\sum_i\mathbf1[y_i\neq c]$ is the modal class.

\textbf{78. Ridge objective strongly convex.} Hessian $2(X^TX+\lambda I)\succeq 2\lambda I\succ0$ for $\lambda>0$.

\textbf{79. Logistic loss decreasing in margin.} $\ell'(z)=-1/(1+e^z)<0$, so larger margin means smaller loss.

\textbf{80. Hinge loss Lipschitz.} $|\max(0,1-a)-\max(0,1-b)|\le |a-b|$; with $a=yw^Tx$, Lipschitz constant in $w$ is $\|x\|$.

\textbf{81. Bayes error $\le 1/2$.} Choosing best class gives conditional error $\min(\eta,1-\eta)\le 1/2$; average over $x$.

\textbf{82. 1NN train error can be 0.} Each point is its own nearest neighbour (or with tie rule/self-inclusion), so memorizes training labels.

\textbf{83. $k$NN prediction is local majority vote.} Classifier picks $\arg\max_y \sum_{i\in N_k(x)}\mathbf1[y_i=y]$.

\textbf{84. Standardization equalizes feature scales.} $x_j'=(x_j-\mu_j)/s_j$ makes each feature mean $0$, variance $1$, preventing distance domination.

\textbf{85. Gini impurity max at uniform binary split.} $G(p)=2p(1-p)$, derivative $2-4p=0$ at $p=1/2$, max $=1/2$.

\textbf{86. Entropy concave.} $-p\log p$ is concave on $(0,1)$; sum preserves concavity, so mixing distributions increases entropy.

\textbf{87. Softmax preserves order of logits.} $a_i>a_j \Rightarrow e^{a_i}>e^{a_j}\Rightarrow q_i>q_j$.

\textbf{88. Logistic MLE = CE minimization.} Negative Bernoulli log-likelihood is exactly cross-entropy, so maximizing likelihood equals minimizing CE.

\textbf{89. KKT complementarity.} For inequality constraint $g_i(w)\le0$, optimality requires $\alpha_i g_i(w)=0$, so active constraints have nonzero multipliers.

\textbf{90. ERM of convex loss is convex.} If each $\ell(w,z_i)$ is convex in $w$, then $L_S(w)=\frac1n\sum_i \ell(w,z_i)$ is convex since nonnegative sums preserve convexity.

\textbf{91. Nonnegative kernel sum is a kernel.} If $K=aK_1+bK_2$ with $a,b\ge0$ and $K_1,K_2$ PSD, then $c^TKc=a\,c^TK_1c+b\,c^TK_2c\ge0$.

\textbf{92. Centering matrix properties.} $C=I-\frac1n\mathbf1\mathbf1^T$ satisfies $C^T=C$ and $C^2=C$, so centering is a symmetric projection.

\textbf{93. Total variance = sum eigenvalues.} For covariance $S=V\Lambda V^T$, $\mathrm{tr}(S)=\mathrm{tr}(\Lambda)=\sum_j\lambda_j$.

\textbf{94. Softmax argmax = logit argmax.} $q_j=e^{a_j}/\sum_k e^{a_k}$; denominator is common and $e^x$ is increasing, so $\arg\max_j q_j=\arg\max_j a_j$.

\end{multicols}
\end{document}
